% Generated from curriculum/manifest.json; do not edit by hand.
\section{两条阅读路线}

\subsection*{应用主线}
先学习数学语言与计算，再进入概率统计、动态模型和数值方法。第 8 章先读有限状态与预测；一般存在性及平方可积理论回读第 3、4 章。第 16、17 章在主线后作为研究实践阅读。

第 1 章 $\rightarrow$ 第 2 章 $\rightarrow$ 第 5 章 $\rightarrow$ 第 6 章 $\rightarrow$ 第 7 章 $\rightarrow$ 第 8 章 $\rightarrow$ 第 9 章 $\rightarrow$ 第 10 章 $\rightarrow$ 第 12 章 $\rightarrow$ 第 13 章 $\rightarrow$ 第 14 章 $\rightarrow$ 第 15 章


\subsection*{理论增强线}
完整学习分析、测度、条件化、概率极限与随机过程，再进入方向模型。

第 1 章 $\rightarrow$ 第 2 章 $\rightarrow$ 第 3 章 $\rightarrow$ 第 4 章 $\rightarrow$ 第 5 章 $\rightarrow$ 第 6 章 $\rightarrow$ 第 7 章 $\rightarrow$ 第 8 章 $\rightarrow$ 第 9 章 $\rightarrow$ 第 11 章

\begin{note}
概率空间、期望交换、条件化与 $L^p$ 论证的完整基础位于第 3、4、8 章。应用主线可按需要回读，理论增强线依次学习这些章节。
\end{note}

\section{应用主线桥接}

\subsection*{概率所需最小测度论桥接}
采用应用主线的读者，可在第 7 章前阅读。需要完整证明时回看第 3 章、第 4 章、第 8 章。

\section{各章所需的基础}

下表列出每章直接用到的前面章节。回顾某章时，也可以沿表查找它所依赖的更早内容。

\begin{longtable}{p{0.07\textwidth}p{0.42\textwidth}p{0.20\textwidth}p{0.20\textwidth}}
\toprule
章 & 学习单元 & 直接先修 & 路线 \\
\midrule
\endfirsthead
\toprule
章 & 学习单元 & 直接先修 & 路线 \\
\midrule
\endhead
1 & 数学语言、证明、反例与量纲 & 无 & 应用主线、理论增强线 \\
2 & 极限、连续、完备性与度量空间 & 第 1 章 & 应用主线、理论增强线 \\
3 & 测度、可测函数与 Lebesgue 积分 & 第 2 章 & 理论增强线 \\
4 & $L^p$ 空间、收敛定理与乘积测度 & 第 3 章 & 理论增强线 \\
5 & 多元微积分与矩阵微分 & 第 1 章、第 2 章 & 应用主线、理论增强线 \\
6 & 线性代数、分解、投影与条件数 & 第 1 章、第 5 章 & 应用主线、理论增强线 \\
7 & 概率空间、随机变量与分布 & 第 1 章 & 应用主线、理论增强线 \\
8 & 条件期望、独立性与概率核 & 第 3 章、第 4 章、第 6 章、第 7 章 & 应用主线、理论增强线 \\
9 & 概率收敛、极限定理与集中不等式 & 第 7 章 & 应用主线、理论增强线 \\
10 & 估计、检验、渐近、回归与多重检验 & 第 6 章、第 9 章 & 应用主线 \\
11 & 随机过程、Markov、鞅、Poisson 与 Brownian & 第 8 章、第 9 章 & 理论增强线 \\
12 & 时间序列与状态空间基础 & 第 10 章 & 应用主线 \\
13 & 凸优化、KKT 与对偶 & 第 5 章、第 6 章 & 应用主线 \\
14 & 浮点数、数值线代与病态问题 & 第 6 章、第 13 章 & 应用主线 \\
15 & Monte Carlo、bootstrap 与方差缩减 & 第 9 章、第 14 章 & 应用主线 \\
16 & 研究有效性、样本外与交易摩擦 & 第 10 章、第 12 章、第 13 章、第 15 章 & 按需 \\
17 & 综合练习：从三笔交易检验研究主张 & 第 16 章 & 按需 \\
\bottomrule
\end{longtable}
